

\documentclass[11pt]{article}
\usepackage[margin=1in]{geometry}
\usepackage[T1]{fontenc}
\usepackage[utf8]{inputenc}
\usepackage{lmodern}
\usepackage{amsmath,amssymb,amsthm}
\usepackage[hidelinks]{hyperref}
\usepackage{enumitem}
\usepackage{microtype}
\usepackage{mathtools}
\usepackage{booktabs}
\usepackage{tabularx}
\usepackage{xurl}
\usepackage[nameinlink,noabbrev]{cleveref}
\usepackage{graphicx}
\usepackage{array}

\newtheorem{theorem}{Theorem}
\newtheorem{definition}{Definition}
\newtheorem{proposition}{Proposition}

\newcommand{\SRP}{\ensuremath{\mathsf{SRP}}}
\newcommand{\tv}{\ensuremath{\mathsf{TV}}}
\newcolumntype{L}[1]{>{\raggedright\arraybackslash}p{#1}}
\newcolumntype{Y}{>{\raggedright\arraybackslash}X}

\title{Auditable Trace Privacy with Abort: Plan-Pinning and Drift Receipts for Side-Channel Resistant Transport}
\author{}
\date{March 24, 2026 (archive v697)}

\begin{document}
\maketitle

\begin{abstract}
Even if a system design satisfies trace-privacy theorems on paper, real deployments drift:
configuration changes, partial rollouts, shared pools, and ``innocent'' performance fixes can silently reintroduce trace leakage.
We propose an operational layer for trace privacy built around \emph{proof-carrying plans}:
(1) every deployment binds itself to a public schedule skeleton (an \SRP plan) via a certificate-like commitment,
(2) records that commitment in a transparency log, and
(3) continuously monitors traces against the plan using \emph{anytime-valid} statistical evidence (e-values / e-processes).
The result is a practical regime of (i) \textbf{fast auditing} (verify what was intended),
(ii) \textbf{split-view risk reduction} (public log consistency plus witnesses/gossip), and
(iii) \textbf{continuous regression detection} (safe under optional stopping).
We provide toy-but-executable tooling: plan certificates, a plan transparency log (PTL), and a trace-audit CLI.
The intended audit stops are likewise small and explicit: plan certificate $\rightarrow$ PTL proof $\rightarrow$ epoch conformance receipt, with evaluation-stealth bias owned by Eval~2 and cross-release packaging deferred to Release~A. The maintained worked side currently records the runtime-conformance stop only as an owner-map / series-spine route: \nolinkurl{example_line_item_owner_map.json} names \nolinkurl{runtime_conformance_packet}, but the current cut does not materialize a \nolinkurl{runtime_conformance_packet} in \nolinkurl{example_uvi.json}, \nolinkurl{eval-runtime-conformance-v1} in the replay-plan catalog, or \nolinkurl{bundle://worked-example/runtime-conformance} in the support map. Later terse papers may use Synthesis~17/Synthesis~55 for that route boundary, while this paper remains the public owner of the local conformance card.\cite{series:workedexample,series:seriesspines}
\end{abstract}

\paragraph{Scope.}
We isolate and formalize \emph{Auditable Trace Privacy with Abort: Plan-Pinning and Drift Receipts for Side-Channel Resistant Transport} as a reusable runtime-conformance contract surface in the anonymity / Anonymous-DHT series.
This paper is the deployment-conformance companion to Operational~A: it owns plan pinning, append-only commitment, and epoch-level drift evidence for a declared runtime plan.
It does not own whether audit traffic is itself distinguishable from live traffic (Eval~2), and it does not own cross-release lineage or signed publication packaging (Release~A).

\paragraph{Units.} Budgets are published in bits (use $\log_2$ and $2^{\epsilon}$); results stated in nats can be converted by dividing by $\ln 2$ (see Synthesis~8).\cite{series:notation}

\paragraph{Imports.}
We treat stop-time padding and the compiler/scheduling machinery as imported from the backbone posts.\cite{series:schedcompiler,series:stoptimeaddendum}
This operational paper specializes that machinery to the retry surface and to schedule-only deployments.
We also treat beacon/VRF binding and append-only log publication as imported primitives (Synthesis~9).\cite{series:transparencyprimitives}
We import interface and threat/window identifiers (ENF-ID and TW-ID) so trace and plan objects can be release-pinned without prose drift.\cite{series:traceifaces,series:threatwindows}

\paragraph{Artifact policy.}
This is a source-only release. Supporting artifacts (logs, traces, scripts, certificates, and verifier outputs) are available on request as a single bundle, committed by ABOM/OINL (Anonymity~C) and governed by the series artifact policy (Synthesis~6).\cite{series:artifactpolicy}

\paragraph{Exports.}
Exports an auditable runtime-conformance contract: (i) a trace schema, (ii) a redaction function defining the public view, and (iii) an epoch conformance receipt certifying a feature-scoped drift budget for a pinned plan.
Deployments additionally export \texttt{plan\_id} and digest-bound \texttt{plan\_spec\_id} (Synthesis~18) together with a plan-certificate inclusion proof in a plan transparency log (PTL).\cite{series:planstateequiv}
This paper deliberately stops before audit-live stealth claims (Eval~2), receipt lineage, and release-ledger comparison; those publication-layer objects are owned elsewhere.

\paragraph{Later-terse-paper citation posture.}
Later terse papers should cite this note only when the live issue is the runtime-conformance surface itself: plan pinning, PTL anchoring, or the epoch conformance receipt. If the live issue is instead the retry-channel theorem root, evaluator detectability / advantage accounting, the mechanical replay checklist for a delivered receipt bundle, the current owner-map / series-spine route for \nolinkurl{runtime_conformance_packet}, theorem-root-to-review routing across that route, or cross-release lineage and package closure, stop at Operational~A / Eval~2 / Synthesis~20 / Synthesis~17--55 / Release~A respectively rather than widening this note into a generic audit-all-purpose companion.\cite{series:operationalA,series:advantagecontracts,series:auditorreplay,series:workedexample,series:seriesspines,series:releaseproperty}

\paragraph{Operational stack position.}
Operational~A supplies the retry-channel theorem root: a pinned public schedule skeleton and a bound of the form $T\delta+\varepsilon(W)$ for the retry trace.
This paper adds only the next deployment-facing stop above that root: a plan certificate, a PTL inclusion proof, and an epoch conformance receipt showing that the live traces still match the pinned plan.
That makes this paper the clean public companion to Operational~A whenever a reader wants deployment-conformance evidence without yet entering the evaluator-stealth or release-lineage layers.
Whether the evaluator's traffic is itself detectably special remains Eval~2, and whether any of these objects become lineaged public release receipts remains Release~A.\cite{series:operationalA,series:advantagecontracts,series:releaseproperty}

\paragraph{Three-stop route above the retry theorem root.}
Later terse papers that need the full operational/public path should cite three stops in order: Operational~A for the declared retry theorem and schedule card, this paper for the deployment-conformance packet, and Release~A for the signed lineage-tagged public release object.
The smallest honest split is: ``Operational~A declared the retry card. Operational~B showed the matching \texttt{plan\_id} / PTL / epoch receipt for that same card. Release~A states whether that audited packet is the same claim surface as its predecessor and what signed digest downstream users may pin.''
This note therefore owns only the middle card: it is not a shadow retry theorem and not a shadow release ledger.\cite{series:operationalA,series:releaseproperty}

\paragraph{Earliest sufficient stop.}
If an auditor can verify (i) the signed plan certificate binding \texttt{plan\_id} / \texttt{plan\_spec\_id} to the deployed build and configuration,
(ii) the PTL inclusion and consistency proofs for that certificate, and (iii) an epoch conformance receipt on the declared public trace features,
then the runtime-conformance question answered by this paper is already settled.
Everything after that---whether audits are themselves detectably special, and whether these objects are rolled into a lineaged public release receipt---belongs to Eval~2 and Release~A respectively.

\section{Introduction}

The companion paper Operational~A (Schedule-only retries for AnonDHT) argues that \SRP wrappers arrest trace leakage from retries, hedges, and negotiation.\cite{series:operationalA}
In Operational~A, the residual risk of a real implementation is captured by an explicit \emph{schedule-only deviation} parameter $\varepsilon(W)$ that adds directly to the trace-privacy bound (while any attempt-level error composes linearly in the public slot budget).
This paper addresses the next question: \emph{how do we keep $\varepsilon(W)$ small over time?} The question here is runtime conformance to a declared plan, not whether the evaluator's traffic is itself stealthy; that audit-awareness problem is separated into Eval~2.

Two failure modes are common:
\begin{itemize}[leftmargin=*]
  \item \textbf{Untracked plan drift.} A deployment silently changes its schedule skeleton (slot counts, lane order, dummy fill policy).
  \item \textbf{In-implementation deviation.} The plan stays constant, but the implementation deviates (bugs, timeouts, pool interactions),
        leaking via micro-structure in the trace.
\end{itemize}

We propose a concrete operational response: \emph{commit, log, and monitor}.
This is inspired by transparency systems in security (Certificate Transparency, CT \cite{rfc6962,rfc9162})
and in the software supply chain (Sigstore/Rekor \cite{newmanSigstore2022,rekorDocs}),
and by safe anytime-valid inference for continuous monitoring \cite{ramdasSAVIStatSci2023,shaferTestingByBetting,ramdasEvaluesEbook2025,lySAVITutorial2025}.

\paragraph{Contributions.}
\begin{itemize}[leftmargin=*]
  \item \textbf{Proof-carrying plans:} a deployment provides a verifiable commitment (a ``plan certificate'') binding code/config to an \SRP schedule skeleton.
  \item \textbf{Plan transparency logs (PTL):} a CT-style append-only log of plan certificates, enabling external auditing and detecting equivocation.
  \item \textbf{Anytime-valid conformance canaries:} e-process based monitoring that remains valid under optional stopping and continuous peeking for drift from a pinned runtime plan.
  \item \textbf{Executable artifacts:} minimal tools to generate/verify plan certificates, append and audit a PTL, and run a trace-audit CLI.
\end{itemize}

\section{Threat Model and Interfaces}
We assume an operator may be honest-but-fallible; an attacker may be external (probing) or internal (a compromised component).
We want to detect:
(i) \emph{misconfiguration} (wrong plan),
(ii) \emph{deviation} (runtime doesn't match the plan), and
(iii) \emph{regressions} over time.

\paragraph{Interface: the plan skeleton.}
A \SRP plan is a public schedule/lane description (slots, lanes, timing class) together with pinned parameters (timeouts, retry budgets).
The plan is not a secret; the point is that \SRP requires that secrets not influence externally visible traces.

\paragraph{Plan identifiers: \texttt{plan\_id} and \texttt{plan\_spec\_id}.}
To align with the receipt schema and drift taxonomy, we use the series' replay-plan identifier pair.\cite{series:receiptitems,series:planstateequiv}
A schedule-only deployment names its plan family by a short \texttt{plan\_id} (e.g., \texttt{plan.srp.traceaudit.v1}) and commits to its semantics by publishing a \texttt{plan\_spec} record.
The \texttt{plan\_spec} includes the public schedule/lane skeleton, pinned parameters (timeouts, retry budgets), and any monitored feature family/thresholds used by this paper's canaries.
Define the digest-bound identifier
\[
\texttt{plan\_spec\_id} := \texttt{sha256}(\texttt{canon}(\texttt{plan\_spec\_core})),
\]
where \texttt{plan\_spec\_core} is the minimal semantic core and \texttt{canon} is a deterministic canonicalization.\cite{series:planstateequiv}
We attach $(\texttt{plan\_id},\texttt{plan\_spec\_id})$ to traces and audit artifacts so ``same plan id, different plan'' cannot hide off-surface.

\paragraph{Binding the interface and the window.}
To keep the operational objects aligned with the series' anonymity claims, the plan object should carry
\texttt{exposure\_nf\_id} (ENF-ID) for the audited trace projection and \texttt{tw\_id} (TW-ID) for the monitoring window.
Including these ids (and their canonical field instantiations) in the hashed plan inputs prevents a deployment from ``keeping the same \texttt{plan\_id}''
while silently changing what it logs or what horizon it claims to monitor.\cite{series:traceifaces,series:threatwindows}


The plan certificate then binds $(\texttt{plan\_id},\texttt{plan\_spec\_id})$ to build metadata (e.g., compiler ID, code hash, and config hash) and is what gets recorded in the PTL.
The explicit domain separator, delimiter, and canonicalization rules are important: auditors should be able to recompute \texttt{plan\_spec\_id} byte-for-byte from published plan inputs.
For JSON-structured inputs, we recommend using a standard canonicalization such as JCS (RFC~8785) to avoid ambiguous encodings; our reference tooling uses a restricted, dependency-free canonicalization suitable for the demo plan types.\cite{rfc8785}

\paragraph{A concrete bridge to $\varepsilon(W)$ in Operational~A: monitoring a feature family.}
Operational~A's end-to-end wrapper deviation $\varepsilon(W)$ is a total-variation quantity over full traces.
Operationally, we pin and monitor a finite family $\mathcal{F}$ of trace features (predicates, counters, lane-use bits) and target the induced deviation (an integral probability metric over $\mathcal{F}$)\cite{sriperumbudurIPM2012}
\[\varepsilon^{\mathcal{F}}(W) := \sup_{\sigma}\;\sup_{s_0,s_1}\;\sup_{f\in\mathcal{F}}\;\big|\mathbb{E}[f(\mathsf{Trace}_{s_0})]-\mathbb{E}[f(\mathsf{Trace}_{s_1})]\big|.\]
If $\mathcal{F}$ is rich enough to separate trace distributions then $\varepsilon^{\mathcal{F}}$ tracks the true $\varepsilon$; otherwise it is a conservative, feature-scoped guarantee.
Crucially, we treat $\mathcal{F}$ and its alerting budget as part of the plan by encoding them in the plan object (e.g., as additional \texttt{pinned\_params} fields). Since \texttt{plan\_spec\_id} hashes both the plan skeleton and \texttt{pinned\_params}, this makes the monitored feature family and its thresholds part of the same identifier that traces and auditors use.

\begin{proposition}[Feature drift vs. total variation (sanity bound)]
Fix a public schedule $\sigma$ and let $P,Q$ be the resulting trace distributions under two secrets.
If every $f\in\mathcal{F}$ takes values in $[-1,1]$, then
\[\sup_{f\in\mathcal{F}}\;\big|\mathbb{E}_P[f]-\mathbb{E}_Q[f]\big| \;\le\; 2\,\tv(P,Q).\]
In particular, any monitored feature drift lower-bounds the true end-to-end deviation: if an audit finds
$\big|\mathbb{E}_P[f]-\mathbb{E}_Q[f]\big|\ge \eta$ for some bounded $f$, then necessarily $\tv(P,Q)\ge \eta/2$.
\end{proposition}

This bound is one-directional: controlling a small feature family does not by itself certify small $\varepsilon(W)$ unless $\mathcal{F}$ is sufficiently expressive.

\paragraph{How to choose $\mathcal{F}$ (practical recipe).}
Pick a small feature family that covers the dominant ways a wrapper can leak.
In practice: (i) \emph{hard conformance predicates} (slot budget exceeded, illegal lane transitions, missing dummy fills),
(ii) \emph{negotiation/fallback bits} (e.g., upgrade lane vs. fallback lane),
(iii) \emph{coarse distributional summaries} (histograms of success slot index, retry counts, abort reason buckets),
and (iv) optionally a \emph{nonparametric IPM feature} such as the maximum mean discrepancy (MMD) computed on low-dimensional trace embeddings.\cite{grettonKernelTwoSample2012,sriperumbudurIPM2012}
We commit $\mathcal{F}$ and its alerting budget in the plan so that auditors can later verify what was monitored. An example MMD-based feature demo is included in \texttt{evidence/mmd\_feature\_demo.py}. More sophisticated anytime-valid two-sample monitoring methods (including sequential and learned-feature variants) can be substituted when appropriate.\cite{pandevaDAVT2024}
For instance, e-values can be constructed for classifier-based two-sample tests (E-C2ST) and for sequential nonparametric two-sample testing by betting, providing additional options when simple handcrafted features are insufficient.\cite{pandevaEC2ST2024,podkopaevSequentialPredictive2023}


\section{Proof-Carrying Plans}
A \emph{plan certificate} is a signed commitment to:
\begin{enumerate}[leftmargin=*]
  \item the plan skeleton (slots, lanes, timing class),
  \item pinned parameters (timeouts, retry budgets),
  \item the code/config identity that enforces the plan (hashes, build metadata).
\end{enumerate}

\paragraph{Signing and verification (practical note).}
In production, the plan certificate should be signed by a deployment identity key (e.g., a build/release service key or an operator key), so third parties can verify that the deployment is actually bound to the published plan.
For interoperability with supply-chain tooling, the certificate can be packaged as a DSSE envelope (a lightweight signature wrapper used broadly for attestations), which reduces reliance on ad-hoc canonicalization rules at the signature layer.\cite{dsseSpec,slsaAttestationModelV12}
Our reference tooling signs the canonicalized certificate fields (excluding the signature itself) with an explicit domain separator using Ed25519, and the PTL treats the serialized certificate as the leaf payload.


\paragraph{Auditor checklist (one page).}
\begin{center}
\begin{tabular}{@{}L{0.22\linewidth}L{0.72\linewidth}@{}}
\toprule
Check & Evidence artifact \\
\midrule
Plan pinned? & Plan certificate (signature) over $(\texttt{plan\_id},\texttt{plan\_spec\_id})$ and policy inputs \\
Discoverable? & PTL inclusion proof for the certificate + latest checkpoint \\
Append-only? & PTL consistency proof between checkpoints (or per-shard checkpoint chain) \\
No split view? & Witness-cosigned checkpoints (or gossip/witness monitors) \\
Runtime matches plan? & Trace contract checks (slot/lane invariants) + canary status \\
\bottomrule
\end{tabular}
\end{center}

This follows the ``proof-carrying'' spirit \cite{neculaPCC97}: a deployment carries a compact object that lets auditors verify key properties quickly.

\paragraph{Trace linkage.}
Each trace record includes the \texttt{plan\_spec\_id} it claims to follow.
The trace-audit CLI first checks that \texttt{plan\_spec\_id} matches the intended plan, then validates per-epoch conformance.
\Cref{fig:cli} shows the tool in action.


\paragraph{Earliest sufficient audit stop.}
If the question is ``was the intended plan pinned and publicly committed?'' an auditor may stop after the signed plan certificate plus PTL inclusion/consistency checks.
If the question is ``did runtime traces drift from the pinned plan this epoch?'' the auditor proceeds to the epoch conformance receipt.
Only if the question is ``how should this epoch compare to other public privacy releases?'' must the audit continue into the Release~A receipt ledger.

\begin{table}[t]
\centering\small
\begin{tabular}{@{}L{0.18\linewidth}L{0.42\linewidth}L{0.28\linewidth}@{}}
\toprule
Field & Minimal public value & Why this row is enough at the conformance stop \\
\midrule
Pinned plan & \texttt{plan\_id=plan.srp.traceaudit.v1}; \texttt{plan\_spec\_id=}digest of canonical plan core & Prevents ``same label, different plan'' drift \\
Claim tuple & \texttt{exposure\_nf\_id=enf.retry.trace.v1}; \texttt{tw\_id=tw.lookup.30d.v1} & Binds the receipt to one declared public trace projection and horizon \\
Epoch scope & \texttt{epoch\_id=2026-W11}; 18{,}432 traces observed & States the exact audit window and sample scope \\
Feature family & Slot-budget checks, lane-transition checks, success-slot histogram, retry-count buckets & Names the monitored proxy for \(\varepsilon^{\mathcal F}(W)\) \\
Contract verdict & No slot/lane violations; no missing dummy fills; \texttt{plan\_spec\_id} matched on every sampled trace & Settles the hard conformance checks before statistics \\
Statistical verdict & Maximum reported e-value 0.61; no canary crossed threshold; declared drift budget for this epoch remained below the public cap & Gives the smallest auditable summary of ``no detected drift this epoch'' \\
Transparency anchor & Certificate inclusion proof at checkpoint \texttt{cp-884}; consistency verified from prior checkpoint & Makes the receipt discoverable and append-only checkable \\
Artifact commitment & Replay bundle digest list committed under ABOM/OINL and available on request & Keeps the public card small without making the private bundle unverifiable \\
\bottomrule
\end{tabular}
\caption{A minimal epoch conformance receipt card. This is the public answer later terse papers can quote directly when the live question is simply whether one declared epoch stayed on one pinned plan.}
\label{tab:minimal-epoch-card}
\end{table}

Table~\ref{tab:minimal-epoch-card} is deliberately non-decorative: it is the smallest public packet that still answers ``was this epoch on the pinned plan?'' without silently widening into evaluator-stealth claims, cross-release lineage, or full replay instructions.

\paragraph{One tiny epoch drill.}
Take the card exactly as written: one pinned plan family, one digest-bound \texttt{plan\_spec\_id}, one epoch \texttt{2026-W11} with 18{,}432 sampled traces, zero hard slot/lane/dummy-fill violations, and a maximum reported e-value of $0.61$ against the declared alert threshold. The public conformance verdict is then intentionally short: \emph{this epoch stayed on the pinned public plan for the declared feature family, and no anytime-valid canary crossed the public drift threshold}. A later terse paper can quote that sentence together with the card rows and stop.

\paragraph{Same card / refreshed-wrapper cutover.}
A later terse paper may keep Table~\ref{tab:minimal-epoch-card} unchanged only when a successor wrapper keeps the same pinned-plan role (same \texttt{plan\_id} family and digest-bound \texttt{plan\_spec\_id} role), the same declared claim tuple (same \texttt{exposure\_nf\_id} and \texttt{tw\_id}), the same monitored feature family, the same public alert rule, and the same transparency-anchor role. In that narrow case the later note may reuse the old conformance sentence verbatim and refresh only wrapper-local lineage labels. If the later note needs a different epoch basis, a different feature family, a different alert threshold rule, or a different PTL/checkpoint role, the shortcut fails closed and this paper must reopen as the owner of the runtime-conformance packet rather than silently treating the old card as ``close enough.''\cite{series:workedexample,series:seriesspines}

\paragraph{One tiny successor-support drill.}
Suppose a successor wrapper changes only the outer release tag while keeping \texttt{plan\_id=plan.srp.traceaudit.v1}, the same canonical-plan digest role for \texttt{plan\_spec\_id}, the same public claim tuple \texttt{(enf.retry.trace.v1, tw.lookup.30d.v1)}, the same slot/lane/retry feature family, the same alert threshold, and the same checkpoint role. Then the later terse paper may say only: \emph{reuse the Operational~B epoch-conformance card unchanged; only the outer wrapper label moved}. That is enough because the conformance packet itself did not change.

\paragraph{One worked-route boundary drill.}
The maintained worked cut names the runtime-conformance stop as an owner-map / series-spine route rather than as a new UVI/support packet. \nolinkurl{example_line_item_owner_map.json} includes \nolinkurl{runtime_conformance_packet} and keeps Operational~A, Operational~B, Boss Fight~A, Eval~2, Synthesis~20, and Release~A in the intended neighboring-owner order. However, \nolinkurl{example_uvi.json} does not contain \nolinkurl{runtime_conformance_packet}; \nolinkurl{example_replay_plans.json} does not contain \nolinkurl{eval-runtime-conformance-v1}; and \nolinkurl{example_support_bundle_map.json} does not contain \nolinkurl{bundle://worked-example/runtime-conformance}. A later terse paper may therefore use the worked owner map and Synthesis~55 only to confirm the route split, while the public conformance claim itself remains the local plan/PTL/epoch receipt card in this paper. Mechanical replay checklists, evaluator detectability transfer, and release lineage remain downstream companions rather than part of this theorem note.\cite{series:workedexample,series:auditorreplay,series:eval1,series:releaseproperty}

\paragraph{Stop / widen rule.}
If a hard conformance predicate fails, if any declared canary crosses its public threshold, or if the same-card / refreshed-wrapper shortcut above is no longer honest, this paper's clean stop is no longer ``green'' and the operator owes an incident/update at the same plan/epoch granularity before anything wider is claimed. Widen to Eval~2 only when the live question is whether the audit traffic or canary regimen is itself detectably special; widen to Synthesis~17 only when the live question is the exact public packet / omitted-support route that carries the epoch receipt onward; widen to Synthesis~55 only when the live question is how this runtime-conformance stop should be cited inside the broader theorem-root-to-review spine; and widen to Release~A only when the live question is cross-release lineage or signed publication closure.\cite{series:advantagecontracts,series:workedexample,series:seriesspines,series:releaseproperty}

\begin{figure}[t]
  \centering
  \fbox{\parbox{0.92\linewidth}{\centering\textbf{Figure omitted in source-only archive.}\\(Plot/diagram and scripts available on request.)}}
  \caption{Trace-audit CLI (from the research archive): traces are validated against a committed plan identifier and monitored for deviations.}
  \label{fig:cli}
\end{figure}

\section{Plan Transparency Log (PTL)}
Plan certificates are more useful if they are publicly \emph{discoverable} and \emph{non-equivocating}.
We follow the standard transparency pattern: an append-only Merkle tree log as in Certificate Transparency (CT) \cite{rfc6962,rfc9162}.
Transparency logs are now largely off the shelf: Certificate Transparency provides a widely deployed
Merkle-tree log interface (checkpoints, inclusion proofs, and consistency proofs), and the same pattern is reused
for non-PKI attestations (e.g., Sigstore-style transparency logs).
Our PTL reuses exactly these interfaces; sharding/tiling and witness cosigning change only how proofs are fetched,
not what an auditor must check.\cite{rfc6962,rfc9162,rekorDocs}

Our toy PTL is a single Merkle tree to keep the paper compact; sharding/tiling does not change the auditor checklist, only how proofs are fetched.\cite{tileBasedLogsTransparencyDev}

\begin{definition}[PTL interface (checkpoint, inclusion, consistency)]
A PTL maintains an append-only Merkle tree over submitted plan certificates (leaf payloads) and exposes:
(i) \emph{signed checkpoints} committing to tree size $n$ and root $r$ (optionally witness-cosigned),
(ii) \emph{inclusion proofs} showing a given leaf is present in the tree at a checkpoint size, and
(iii) \emph{consistency proofs} showing that a later checkpoint extends an earlier checkpoint without rewriting prior leaves.
\end{definition}

\begin{proposition}[Append-only auditability (informal)]
If a plan certificate verifies and has a valid inclusion proof at checkpoint $\mathsf{cp}_2$, and auditors also verify a consistency proof from an earlier checkpoint $\mathsf{cp}_1$ to $\mathsf{cp}_2$, then the log view at $\mathsf{cp}_2$ is an append-only extension of the view at $\mathsf{cp}_1$ (no fork that rewrites history without being detected).
\end{proposition}


\paragraph{What auditors check.}
Given a plan certificate with \texttt{plan\_spec\_id}, auditors verify:
(i) a valid signature,
(ii) an inclusion proof for the certificate in the PTL, and
(iii) a consistency proof between PTL checkpoints.
This detects silent plan drift and supports third-party review.
\Cref{fig:ptl} shows a toy-but-executable PTL.

\begin{figure}[t]
  \centering
  \fbox{\parbox{0.92\linewidth}{\centering\textbf{Figure omitted in source-only archive.}\\(Plot/diagram and scripts available on request.)}}
  \caption{Toy plan transparency log (PTL) demo: append-only Merkle structure enables inclusion/consistency proofs, as in CT.}
  \label{fig:ptl}
\end{figure}

\paragraph{Witnessing / cosigning.}
Transparency logs can still equivocate if clients do not gossip.
Witness cosigning strengthens logs by proactively recording checkpoints with independent witnesses, as advocated by CoSi \cite{sytaCoSi2016} and formalized in modern tlog cosignature/witness interfaces.\cite{c2spTlogCosignature,c2spTlogWitness}
Operational efforts package witnesses as plug-and-play components (e.g., ArmoredWitness) and publish configuration guidance for deploying witness sets; some operators also publish concrete witness endpoints and conformance statements.\cite{armoredWitnessGithub,trustfabricWitnesses2025}
\Cref{fig:witness} illustrates a witness-cosigned checkpoint flow.

\begin{figure}[t]
  \centering
  \fbox{\parbox{0.92\linewidth}{\centering\textbf{Figure omitted in source-only archive.}\\(Plot/diagram and scripts available on request.)}}
  \caption{Witness cosigning for PTL checkpoints (toy demo): witness-signed checkpoints reduce the risk of undetected equivocation.}
  \label{fig:witness}
\end{figure}

\section{Anytime-Valid Canary Monitoring}
Auditing detects \emph{what was committed}; monitoring detects \emph{what is happening}.

\paragraph{Why anytime-valid?}
Operational monitoring involves continuous peeking (dashboards, alerts).
Classical p-values inflate false positive rates under optional stopping.
E-values and e-processes are designed for time-uniform validity \cite{ramdasSAVIStatSci2023,shaferTestingByBetting,ramdasEvaluesEbook2025}.
\Cref{fig:optional} visualizes the difference.

\paragraph{Where do the samples come from? (conformance canaries)}
To monitor a feature family $\mathcal{F}$ in a way that matches Operational~A's ``two secrets'' framing, operators can run \emph{randomized canary sessions}.
In each epoch (or test window), choose a canary label $C_t\in\{0,1\}$ using public randomness, execute an otherwise plan-identical canary request from a small set of pinned canary inputs, and log the plan-defined feature vector $X_t$ of the observed trace.
Under end-to-end \SRP, $X_t$ should be independent of $C_t$ (conditioned on the public plan), so any detected dependence is actionable evidence of drift and directly lower-bounds the feature-scoped deviation $\varepsilon^{\mathcal{F}}(W)$.
Our lane-switch regression toy below can be read as this setting when $X_t$ is the single binary feature ``fallback lane used.''

\paragraph{Note on thresholds.}
The default rejection threshold $1/\alpha$ (from Markov/Ville) is intentionally conservative; under additional, empirically checkable distributional assumptions on e-values, the threshold can sometimes be lowered while retaining Type~I control, improving detection power.\cite{blierWongWangImprovedThresholds2026}

\begin{figure}[t]
  \centering
  \fbox{\parbox{0.92\linewidth}{\centering\textbf{Figure omitted in source-only archive.}\\(Plot/diagram and scripts available on request.)}}
  \caption{Optional stopping intuition (from the research archive): classical p-values break under repeated peeking; e-values/e-processes remain valid at all stopping times.}
  \label{fig:optional}
\end{figure}

\begin{theorem}[Time-uniform false-alarm control (Ville/Doob maximal inequality, informal)]
Let $(E_t)_{t\ge 0}$ be a nonnegative \emph{e-process} under $H_0$, i.e., a nonnegative supermartingale with $E_0=1$.
Define a stopping rule $\tau := \inf\{t : E_t \ge 1/\alpha\}$.
Then $\Pr_{H_0}[\tau < \infty] \le \alpha$.\cite{villeCollectif1939,ramdasEvaluesEbook2025,howardConfidenceSequencesAOS2021,rufCompositeVille2023,koolenVille2026}
\end{theorem}
\noindent\emph{Sanity check.} The accompanying evidence scripts include an empirical $H_0$ false-alarm check over finite horizons (not plotted here).


\paragraph{A concrete regression: lane-switch / downgrade.}
A common \SRP failure mode is negotiation fallback becoming data-dependent (or plan-inconsistent),
turning ``fallback used'' into an amplifiable bit.
We include a minimal synthetic experiment: a single binary feature per epoch (fallback lane used) with baseline probability $p_0$.
Under regression, the probability increases by $\Delta$.
Using a likelihood-ratio e-process, we estimate detection times vs. $\Delta$ while preserving time-uniform false-alarm control.
For a single binary feature $X_t\in\{0,1\}$ with null rate $p_0$ and an alternative rate $p_1>p_0$, the canonical per-epoch likelihood ratio is
\[L_t(p_1)=\Big(\frac{p_1}{p_0}\Big)^{X_t}\Big(\frac{1-p_1}{1-p_0}\Big)^{1-X_t},\qquad E_T(p_1)=\prod_{t=1}^{T} L_t(p_1),\]
which is an e-process under $H_0$; in practice we use a simple mixture over a small grid of $p_1$ values to avoid picking a single alternative.
\Cref{fig:laneswitch} also overlays a back-of-the-envelope approximation
$\log_2(1/\alpha)/\mathrm{KL}(p_1\|p_0)$.

\begin{figure}[t]
  \centering
  \fbox{\parbox{0.92\linewidth}{\centering\textbf{Figure omitted in source-only archive.}\\(Plot/diagram and scripts available on request.)}}
  \caption{Synthetic lane-switch regression detection (script and CSV in \texttt{evidence/}): larger downgrade shifts are detected faster; the $\log_2(1/\alpha)/\mathrm{KL}$ curve provides a useful sizing heuristic.}
  \label{fig:laneswitch}
\end{figure}

\paragraph{Multi-feature monitoring (omitted for space).}
The research archive includes a multifeature scaling experiment showing how combining trace features can improve detection speed while preserving time-uniform guarantees.


\section{Sizing (brief)}
There is an inherent race: adversaries can exploit deviations quickly, while auditors need enough samples to detect them.
The research archive includes an ``attack-vs-audit'' sizing worksheet and plot; we omit the figure here for space.


\section{Operational playbook (condensed)}
A practical playbook:
\begin{enumerate}[leftmargin=*]
  \item \textbf{Pin} a plan skeleton and produce a plan certificate at build/rollout time.
  \item \textbf{Append} the certificate to the PTL (obtain inclusion proof; record checkpoint).
  \item \textbf{Monitor} traces against the plan (contract checks) and run anytime-valid canaries for drift.
  \item \textbf{Respond} to alarms: freeze rollouts, compare deployed \texttt{plan\_spec\_id} to expected, and publish incident notes.
\end{enumerate}

\begin{figure}[t]
  \centering
  \fbox{\parbox{0.92\linewidth}{\centering\textbf{Figure omitted in source-only archive.}\\(Plot/diagram and scripts available on request.)}}
  \caption{Audit playbook flowchart (from the research archive).}
  \label{fig:playbook}
\end{figure}

\section{Related Work (selected)}
Certificate Transparency \cite{rfc6962,rfc9162} is the canonical transparency-log design for PKI.
Key transparency systems (e.g., CONIKS and Google Key Transparency) apply similar append-only log and monitoring ideas to public key directories and account keys.\cite{coniks2015,keyTransparencyGoogle}
Studies of CT in practice emphasize the importance of monitor diversity, completeness, and anti-equivocation techniques \cite{liCTInWild2019,sunCTMonitorsNDSS2024,dahlbergAggCTGossip2019}.
Sigstore applies similar ideas to software artifacts via a public log (Rekor) \cite{newmanSigstore2022,rekorDocs}. Related transparency frameworks are being explored in other domains (e.g., confidential computing transparency for TEEs) \cite{cctTransparency2024}.
Supply-chain attestation frameworks (e.g., in-toto and SLSA) and emerging IETF work on SCITT likewise center on signed statements plus transparency/receipts (including COSE Receipts with Merkle inclusion/consistency proofs), making PTL-style logs a natural fit for plan certificates.\cite{torresAriasInToto2019,slsaSpecV12,scittArchitectureDraft2025,scittWGAbout2026,scittReceiptsCcfProfile2026,coseMerkleTreeProofs2025}
Witness cosigning (CoSi) strengthens transparency services by proactively preventing undetected equivocation \cite{sytaCoSi2016}.
For the monitoring layer, safe anytime-valid inference (e-values/e-processes) enables optional-stopping-robust alerting \cite{ramdasSAVIStatSci2023,shaferTestingByBetting,ramdasEvaluesEbook2025}.

\section{Limitations (brief)}
\label{sec:c2limitations}
This paper focuses on making trace-privacy \emph{auditable} and \emph{monitorable}, not on hiding the fact that a plan exists.
Because plans and checkpoints are public, they can reveal operational metadata (e.g., which schedule families a service supports); if that is sensitive, deployments can aggregate or delay publication, or publish only a coarse plan family while keeping fine-grained parameters internal.
Finally, the monitoring layer is only as good as the committed feature family $\mathcal{F}$: if drift manifests outside $\mathcal{F}$ it may go undetected until $\mathcal{F}$ is expanded or the plan is revised.

\section{Conclusion}
Trace privacy needs governance.
Proof-carrying plans, transparency logs, and anytime-valid monitoring provide a practical path
from a paper design to an auditable, regression-resistant deployment discipline.
The clean public stop of this paper is still pre-release: PTL-backed plan pinning plus an epoch conformance receipt, not a lineaged publication receipt.

\section{Takeaways}
\begin{itemize}[leftmargin=*]
  \item \textbf{Pin the plan:} a deployment exports a \texttt{plan\_spec\_id} and a signed plan certificate binding \texttt{plan\_spec\_id} to build/config hashes; the certificate is appended to a PTL.
  \item \textbf{Make drift checkable:} traces carry \texttt{plan\_spec\_id}; auditors verify conformance and publish an inclusion proof plus an epoch conformance receipt.
  \item \textbf{Monitor under peeking:} e-process/e-value receipts bound feature-level drift $\varepsilon^{\mathcal{F}}(W)$ under optional stopping, giving an operational proxy for Operational~A's schedule-only deviation.
  \item \textbf{Artifacts (on request):} reference tooling (\path{tools/}) and demos (\path{evidence/}) are available on request under the series artifact policy.\cite{series:artifactpolicy}
  \item \textbf{Keep the seam clean:} this paper owns runtime conformance to a pinned plan; Eval~2 owns audit-live detectability bias, and Release~A owns cross-release lineage and publication packaging.
\end{itemize}

\begin{thebibliography}{99}\small

\bibitem{neculaPCC97}
George C. Necula.
\newblock Proof-Carrying Code.
\newblock Proceedings of the 24th ACM SIGPLAN--SIGACT Symposium on Principles of Programming Languages (POPL), 106--119, 1997.
\newblock DOI: 10.1145/263699.263712.

\bibitem{rfc6962}
Ben Laurie and Adam Langley and Emilia Kasper.
\newblock Certificate Transparency.
\newblock IETF, 2013.
\newblock DOI: 10.17487/RFC6962.
\newblock \url{https://www.rfc-editor.org/rfc/rfc6962.html}

\bibitem{rfc9162}
Ben Laurie and Emilia Messeri and Rob Stradling.
\newblock Certificate Transparency Version 2.0.
\newblock IETF, 2021.
\newblock DOI: 10.17487/RFC9162.
\newblock \url{https://www.rfc-editor.org/rfc/rfc9162.html}

\bibitem{trillianGeneralTransparency}
{Google Trillian Project}.
\newblock Trillian: General Transparency (Log Mode Overview).
\newblock Project documentation, 2026.
\newblock States Trillian is in maintenance mode and recommends Tessera for new tiled transparency logs. Accessed 2026-02-22.
\newblock \url{https://google.github.io/trillian/}

\bibitem{tesseraGithub}
{transparency-dev}.
\newblock Tessera: Go library for building tile-based transparency logs (tlogs).
\newblock GitHub repository, 2024.
\newblock Describes Tessera as a tile-based successor to Trillian v1 and includes witness and Static CT API support. Accessed 2026-02-21.
\newblock \url{https://github.com/transparency-dev/tessera}

\bibitem{tileBasedLogsTransparencyDev}
{transparency.dev}.
\newblock Tile-Based Transparency Logs.
\newblock Technical article, 2024.
\newblock Overview of tiled log APIs (checkpoints, tiles, and proof distribution) and operational motivations. Accessed 2026-02-21.
\newblock \url{https://transparency.dev/articles/tile-based-logs/}

\bibitem{tesseraCTGithub}
{transparency-dev}.
\newblock TesseraCT (\texttt{tesseract}): Static CT API server built on Tessera.
\newblock GitHub repository, 2025.
\newblock Implements the Static CT API on top of Tessera for production-grade CT logs. Accessed 2026-02-21.
\newblock \url{https://github.com/transparency-dev/tesseract}

\bibitem{armoredWitnessGithub}
{transparency.dev}.
\newblock ArmoredWitness: A plug-and-play witness solution.
\newblock GitHub repository, 2024.
\newblock Project intended to kick-start a cross-ecosystem witness network and reduce split-view risk. Accessed 2026-02-21.
\newblock \url{https://github.com/transparency-dev/armored-witness}

\bibitem{newmanSigstore2022}
Zachary Newman and others.
\newblock Sigstore: Software Signing for Everybody.
\newblock ACM Conference on Computer and Communications Security (CCS), 2022.
\newblock DOI: 10.1145/3548606.3560596.

\bibitem{rekorDocs}
{Sigstore Project}.
\newblock Rekor Transparency Log: Overview.
\newblock Documentation, 2026.
\newblock \url{https://docs.sigstore.dev/logging/overview/}

\bibitem{torresAriasInToto2019}
Torres-Arias, Santiago and others.
\newblock in-toto: Providing farm-to-table guarantees for bits and bytes.
\newblock 28th {USENIX} Security Symposium ({USENIX} Security 19), 2019.
\newblock \url{https://www.usenix.org/system/files/sec19-torres-arias.pdf}

\bibitem{dsseSpec}
{Secure Systems Lab}.
\newblock DSSE: Dead Simple Signing Envelope (specification).
\newblock GitHub repository (spec), 2026.
\newblock Defines a lightweight envelope format for signing arbitrary payloads; used by multiple supply-chain attestation systems. Accessed 2026-02-22.
\newblock \url{https://github.com/secure-systems-lab/dsse}

\bibitem{slsaSpecV12}
{OpenSSF / SLSA}.
\newblock SLSA Specification (v1.2).
\newblock Specification, 2026.
\newblock Approved specification; accessed 2026-02-22.
\newblock \url{https://slsa.dev/spec/v1.2/}

\bibitem{slsaAttestationModelV12}
{OpenSSF / SLSA}.
\newblock SLSA v1.2: Software Attestations (recommended envelope: DSSE).
\newblock Specification page, 2026.
\newblock Defines a recommended suite of attestation building blocks and lists DSSE as the recommended envelope. Accessed 2026-02-22.
\newblock \url{https://slsa.dev/spec/v1.2/attestation-model}

\bibitem{scittArchitectureDraft2025}
Birkholz, Henk and others.
\newblock An Architecture for Trustworthy and Transparent Digital Supply Chains (SCITT).
\newblock Internet-Draft, \texttt{draft-ietf-scitt-architecture-22}, 2025.
\newblock Expires 13 April 2026; accessed 2026-02-22.
\newblock \url{https://datatracker.ietf.org/doc/draft-ietf-scitt-architecture/}

\bibitem{rekorV2ga2025}
{Sigstore Project}.
\newblock Rekor v2 GA: Cheaper to run, simpler to maintain.
\newblock Blog post, 2025.
\newblock \url{https://blog.sigstore.dev/rekor-v2-ga/}

\bibitem{rfc8785}
Rundgren, Anders.
\newblock JSON Canonicalization Scheme (JCS).
\newblock IETF, 2020.
\newblock DOI: 10.17487/RFC8785.
\newblock \url{https://www.rfc-editor.org/rfc/rfc8785}

\bibitem{sriperumbudurIPM2012}
Sriperumbudur, Bharath K. and Fukumizu, Kenji and Gretton, Arthur and Sch{"o}lkopf, Bernhard and Lanckriet, Gert R. G..
\newblock On the Empirical Estimation of Integral Probability Metrics.
\newblock Electronic Journal of Statistics, 6, 1550--1599, 2012.
\newblock DOI: 10.1214/12-EJS722.
\newblock \url{https://projecteuclid.org/journals/electronic-journal-of-statistics/volume-6/issue-none/On-the-empirical-estimation-of-integral-probability-metrics/10.1214/12-EJS722.full}

\bibitem{sytaCoSi2016}
Ewa Syta and Iulia Tamas and Dylan Visher and David Isaac Wolinsky and Philipp Jovanovic and Linus Gasser and Nicolas Gailly and Ismail Khoffi and Bryan Ford.
\newblock Keeping Authorities ``Honest or Bust'' with Decentralized Witness Cosigning.
\newblock IEEE Symposium on Security and Privacy (S\&P), 2016.
\newblock DOI: 10.1109/SP.2016.38.
\newblock \url{https://www.davidwolinsky.com/site/files/syta16sp.pdf}

\bibitem{ramdasSAVIStatSci2023}
Aaditya Ramdas and Peter Gr\"unwald and Vladimir Vovk and Glenn Shafer.
\newblock Game-Theoretic Statistics and Safe Anytime-Valid Inference.
\newblock Statistical Science, 38(4), 2023.
\newblock DOI: 10.1214/23-STS894.
\newblock \url{https://projecteuclid.org/journals/statistical-science/volume-38/issue-4/Game-Theoretic-Statistics-and-Safe-Anytime-Valid-Inference/10.1214/23-STS894.full}

\bibitem{villeCollectif1939}
Ville, Jean.
\newblock \'{E}tude critique de la notion de collectif.
\newblock 1939.
\newblock Introduces Ville's martingale inequality underlying time-uniform bounds; scanned thesis/book hosted by Numdam. Accessed 2026-02-22.
\newblock \url{https://www.numdam.org/item/THESE_1939__218__1_0.pdf}

\bibitem{howardConfidenceSequencesAOS2021}
Howard, Steven R. and Ramdas, Aaditya and McAuliffe, Jon and Sekhon, Jasjeet.
\newblock Time-uniform, nonparametric, nonasymptotic confidence sequences.
\newblock The Annals of Statistics, 49(2), 2021.
\newblock DOI: 10.1214/20-AOS1991.
\newblock \url{https://projecteuclid.org/journals/annals-of-statistics/volume-49/issue-2/Time-uniform-nonparametric-nonasymptotic-confidence-sequences/10.1214/20-AOS1991.full}

\bibitem{letsencrypt6962eol}
{Let's Encrypt}.
\newblock End of Life Plan for RFC 6962 Certificate Transparency Logs.
\newblock Blog post, 2025.
\newblock \url{https://letsencrypt.org/2025/08/14/rfc-6962-logs-eol}

\bibitem{shaferTestingByBetting}
Shafer, Glenn.
\newblock Testing by Betting: A Strategy for Statistical and Scientific Communication.
\newblock Journal of the Royal Statistical Society: Series A, 2021.
\newblock \url{https://glennshafer.com/assets/downloads/articles/article104_jrss_shafer-testingbybetting-with-discussion-response.pdf}

\bibitem{ramdasEvaluesEbook2025}
Ramdas, Aaditya and others.
\newblock Hypothesis Testing with E-values.
\newblock Monograph (draft), 2025.
\newblock \url{https://www.stat.cmu.edu/~aramdas/ebook-final.pdf}

\bibitem{rufCompositeVille2023}
Ruf, Johannes and Larsson, Martin and Koolen, Wouter M. and Ramdas, Aaditya.
\newblock A Composite Generalization of Ville's Martingale Theorem Using e-Processes.
\newblock Electronic Journal of Probability, 28, 1--21, 2023.
\newblock DOI: 10.1214/23-EJP1019.
\newblock \url{https://projecteuclid.org/journals/electronic-journal-of-probability/volume-28/issue-none/A-composite-generalization-of-Villes-martingale-theorem-using-e-processes/10.1214/23-EJP1019.full}

\bibitem{blierWongWangImprovedThresholds2026}
Blier-Wong, Christopher and Wang, Ruodu.
\newblock Improved Thresholds for E-values.
\newblock Preprint, 2026.
\newblock Discusses improved rejection thresholds for e-tests/e-processes under additional distributional assumptions; dated Feb 7, 2026. Accessed 2026-02-22.
\newblock \url{https://www.math.uwaterloo.ca/~wang/papers/2026BlierWong-Wang-AOS.pdf}

\bibitem{dahlbergAggCTGossip2019}
Rasmus Dahlberg and Tobias Pulls and Roel Peeters.
\newblock Aggregation-Based Certificate Transparency Gossip.
\newblock arXiv:1806.08817, 2018.
\newblock \url{https://arxiv.org/abs/1806.08817}

\bibitem{liCTInWild2019}
Bingyu Li and Chuan Qin and Jianjun Chen and others.
\newblock Certificate Transparency in the Wild.
\newblock ACM Conference on Computer and Communications Security (CCS), 2019.
\newblock \url{https://www.inforsec.org/wp/wp-content/uploads/2020/04/CCS2019-Li.pdf}

\bibitem{sunCTMonitorsNDSS2024}
Aozhuo Sun and Jinqiang Lin and Wei Wang and others.
\newblock The Public Inspections on Third-party Monitors.
\newblock Network and Distributed System Security Symposium (NDSS), 2024.
\newblock \url{https://www.ndss-symposium.org/wp-content/uploads/2024-834-paper.pdf}

\bibitem{lySAVITutorial2025}
Alexander Ly and Udo Boehm and Peter Gr\"unwald and Aaditya Ramdas and Don van Ravenzwaaij.
\newblock A Tutorial on Safe Anytime-Valid Inference: Practical Maximally Flexible Sampling Designs for Experiments Based on e-Values.
\newblock Preprint, 2025.
\newblock \url{https://www.alexander-ly.com/wp-content/uploads/2025/08/saviTutorial.pdf}

\bibitem{pandevaDAVT2024}
Pandeva, Teodora and Forr{\'e}, Patrick and Ramdas, Aaditya and Shekhar, Shubhanshu.
\newblock Deep Anytime-Valid Hypothesis Testing.
\newblock Proceedings of The 27th International Conference on Artificial Intelligence and Statistics (AISTATS), 238, 622--630, 2024.
\newblock \url{https://proceedings.mlr.press/v238/pandeva24a.html}

\bibitem{pandevaEC2ST2024}
Pandeva, Teodora and Bakker, Tim and Naesseth, Christian A. and Forr{\'e}, Patrick.
\newblock E-Valuating Classifier Two-Sample Tests.
\newblock Transactions on Machine Learning Research (TMLR), 2024.
\newblock Introduces e-values for classifier two-sample tests (E-C2ST) suitable for anytime-valid sequential testing.
\newblock \url{https://openreview.net/pdf?id=dwFRov8xhr}

\bibitem{podkopaevSequentialPredictive2023}
Podkopaev, Aleksandr and Ramdas, Aaditya.
\newblock Sequential Predictive Two-Sample and Independence Testing.
\newblock Advances in Neural Information Processing Systems (NeurIPS), 2023.
\newblock Nonparametric sequential two-sample testing by betting with anytime-valid guarantees.
\newblock \url{https://proceedings.neurips.cc/paper_files/paper/2023/file/a72b207734d6112f6b47447e46be40e9-Paper-Conference.pdf}

\bibitem{koolenVille2026}
Koolen, Wouter M. and P{'e}rez-Ortiz, Muriel F. and Lardy, Tyron.
\newblock A Generalisation of Ville's Inequality to Monotonic Lower Bounds and Thresholds.
\newblock Statistics \& Probability Letters, 229, 110577, 2026.
\newblock DOI: 10.1016/j.spl.2025.110577.
\newblock \url{https://ir.cwi.nl/pub/35761/35761.pdf}

\bibitem{grettonKernelTwoSample2012}
Gretton, Arthur and Borgwardt, Karsten M. and Rasch, Malte J. and Sch{"o}lkopf, Bernhard and Smola, Alexander J..
\newblock A Kernel Two-Sample Test.
\newblock Journal of Machine Learning Research, 13, 723--773, 2012.
\newblock \url{https://www.jmlr.org/papers/volume13/gretton12a/gretton12a.pdf}

\bibitem{cctTransparency2024}
Kocao\u{g}ullar, Ceren and Marjanov, Tina and Petrov, Ivan and Laurie, Ben and Cutter, Al and Kern, Christoph and Hutchings, Alice and Beresford, Alastair R..
\newblock Confidential Computing Transparency.
\newblock 2024.
\newblock arXiv:2409.03720
\newblock \url{https://arxiv.org/abs/2409.03720}

\bibitem{coniks2015}
Melara, Marc S. and Blankstein, Aaron and Bonneau, Joseph and Felten, Edward W. and Freedman, Michael J..
\newblock CONIKS: Bringing Key Transparency to End Users.
\newblock Proceedings of the 24th USENIX Security Symposium, 2015.
\newblock \url{https://www.usenix.org/system/files/conference/usenixsecurity15/sec15-paper-melara.pdf}

\bibitem{keyTransparencyGoogle}
{Google}.
\newblock Key Transparency (open-source implementation).
\newblock GitHub repository, 2024.
\newblock \url{https://github.com/google/keytransparency}

\bibitem{c2spTlogTiles}
{C2SP}.
\newblock Tiled Transparency Logs (tlog-tiles).
\newblock Community specification, 2025.
\newblock Specifies an HTTP API for fetching checkpoints, tiles, and entries of a tiled transparency log. Accessed 2026-02-22.
\newblock \url{https://c2sp.org/tlog-tiles}

\bibitem{c2spStaticCtApi}
{C2SP}.
\newblock Static Certificate Transparency API (static-ct-api).
\newblock Community specification, 2025.
\newblock Specifies the public endpoints of the Static CT API used by next-generation CT logs. Accessed 2026-02-22.
\newblock \url{https://c2sp.org/static-ct-api}

\bibitem{c2spTlogCheckpoint}
\newblock C2SP: tlog-checkpoint.
\newblock Community Cryptography Specification Project.
\newblock Accessed 2026-02-22
\newblock \url{https://c2sp.org/tlog-checkpoint}

\bibitem{c2spTlogProof}
\newblock C2SP: tlog-proof.
\newblock Community Cryptography Specification Project.
\newblock Accessed 2026-02-22
\newblock \url{https://c2sp.org/tlog-proof}

\bibitem{c2spTlogCosignature}
\newblock C2SP: tlog-cosignature.
\newblock Community Cryptography Specification Project.
\newblock Accessed 2026-02-22
\newblock \url{https://c2sp.org/tlog-cosignature}

\bibitem{c2spTlogWitness}
\newblock C2SP: tlog-witness.
\newblock Community Cryptography Specification Project.
\newblock Accessed 2026-02-22
\newblock \url{https://c2sp.org/tlog-witness}

\bibitem{chromeCTPolicyPage2026}
{Google Chrome}.
\newblock Chrome Certificate Transparency Policy.
\newblock Web page, 2026.
\newblock Accessed 2026-02-22; includes the transition rule that before April 15, 2026 at least one SCT must be from an RFC6962-compliant log.
\newblock \url{https://googlechrome.github.io/CertificateTransparency/ct_policy.html}

\bibitem{letsencryptCTLogsDoc2026}
{Let's Encrypt}.
\newblock Certificate Transparency (CT) Logs.
\newblock Documentation page, 2026.
\newblock Summarizes RFC 6962 log EOL status and shutdown timeline; updated 2026-01-04.
\newblock \url{https://letsencrypt.org/docs/ct-logs/}

\bibitem{valsordaKeyserverTlog2025}
Filippo Valsorda.
\newblock Building a Transparent Keyserver.
\newblock Technical article, 2025.
\newblock Step-by-step integration of a tiled transparency log (Tessera) and witness cosigning (Witness Network) into a centralized service. Accessed 2026-02-22.
\newblock \url{https://words.filippo.io/keyserver-tlog/}

\bibitem{scittWGAbout2026}
{IETF SCITT Working Group}.
\newblock Supply Chain Integrity, Transparency, and Trust ({SCITT}) Working Group.
\newblock IETF Datatracker (group page), 2026.
\newblock Accessed 2026-02-22.
\newblock \url{https://datatracker.ietf.org/group/scitt/about/}

\bibitem{scittReceiptsCcfProfile2026}
{IETF SCITT Working Group}.
\newblock COSE Receipts with CCF.
\newblock Internet-Draft, \texttt{draft-ietf-scitt-receipts-ccf-profile-00}, 2026.
\newblock Accessed 2026-02-22.
\newblock \url{https://datatracker.ietf.org/doc/draft-ietf-scitt-receipts-ccf-profile/}

\bibitem{trustfabricWitnesses2025}
{Transparency.dev (TrustFabric team)}.
\newblock TrustFabric Witnesses.
\newblock Project page, 2025.
\newblock Lists witnesses and notes conformance to C2SP tlog-witness and related specs. Accessed 2026-02-22.
\newblock \url{https://transparency.dev/witnesses/}

\bibitem{plantsMerkleTreeCerts2026}
Benjamin, David and others.
\newblock Merkle Tree Certificates.
\newblock Internet-Draft, \texttt{draft-ietf-plants-merkle-tree-certs-01}, 2026.
\newblock IETF PLaNTS draft extending RFC 9162 with subtree proof objects and cosigner roles. Accessed 2026-02-22.
\newblock \url{https://datatracker.ietf.org/doc/draft-ietf-plants-merkle-tree-certs/}

\bibitem{coseMerkleTreeProofs2025}
Orie Steele and Henk Birkholz and Amaury Delignat-Lavaud and Cedric Fournet.
\newblock COSE (CBOR Object Signing and Encryption) Receipts.
\newblock Internet-Draft, \texttt{draft-ietf-cose-merkle-tree-proofs-18}, 2025.
\newblock Accessed 2026-02-22.
\newblock \url{https://datatracker.ietf.org/doc/draft-ietf-cose-merkle-tree-proofs/}


\bibitem{series:transparencyprimitives}
Anonymous.
\newblock Transparency Primitives for Proof-Carrying Anonymity Claims: Beacons, VRFs, and Append-Only Logs.
\newblock Series synthesis note (Synthesis~9), The-One-True-Archive (2026).

\bibitem{series:traceifaces}
Anonymous.
\newblock Trace Interfaces for Anonymous DHTs: Observation Tiers, Committee-Contact Traces, and Exposure Normal Forms.
\newblock Series synthesis note (Synthesis~3), 2026. (This archive.)

\bibitem{series:threatwindows}
Anonymous.
\newblock Threat Models and Repetition Windows for Anonymous DHTs: A Short Declaration Checklist for Referee-Grade Budget Claims.
\newblock Series synthesis note (Synthesis~4), 2026. (This archive.)

\bibitem{series:receiptitems}
Anonymous.
\newblock Receipt Line Items for Anonymous-DHT Privacy Claims: Minimal Tuple Schema and Drift Semantics.
\newblock Series synthesis note (Synthesis~12), 2026. (This archive.)

\bibitem{series:planstateequiv}
Anonymous.
\newblock Digest-Bound Replay Plans and State Declarations: Equivalence, Minimality, and Change Control.
\newblock Series synthesis note (Synthesis~18), 2026. (This archive.)

\bibitem{series:artifactpolicy}
Anonymous.
\newblock Artifact policy and supporting materials for the anonymity/AnonDHT series.
\newblock 2026.
\newblock See synthesis paper: Artifact policy (this archive). Artifacts, tooling, certificates, and scripts are available on request as a single archive.



\bibitem{series:operationalA}
Anonymous DHT Series (draft).
\emph{Schedule-Only Retries for AnonDHT: The SRP Wrapper Contract}.
In this archive: Operational Series, Paper~A.

\bibitem{series:schedcompiler}
Anonymous DHT Series (published).
\emph{Scheduling and Compiler Techniques for Metadata-Hiding Overlay Lookups}.
\texttt{[[2026.01.22 - Mathematics: Scheduling and Compiler Techniques for Metadata-Hiding Overlay Lookups]]}.

\bibitem{series:stoptimeaddendum}
Anonymous DHT Series (published).
\emph{Stop-Time Padding Addendum, Max-Mixture Separations and Tail-Sign Deadline Normal Forms}.
\texttt{[[2026.01.26 - Mathematics: Stop-Time Padding Addendum, Max-Mixture Separations and Tail-Sign Deadline Normal Forms]]}.


\bibitem{series:notation}
Anonymous.
\newblock Notation and Minimal Model for the One-True-Archive.
\newblock Series synthesis note (Synthesis~8), 2026. (This archive.)


\bibitem{series:advantagecontracts}
Anonymous.
\newblock Advantage Contracts for Anonymous DHT Evaluations: Stealth Audits, Leakage Discovery at Scale, and Practical Padding Filters.
\newblock Evaluation series Paper~2, 2026. (This archive.)

\bibitem{series:auditorreplay}
Anonymous.
\newblock Auditor Replay for Anonymous-DHT Receipt Line Items: A Mechanical Checklist for Digest-Bound Claims.
\newblock Series synthesis note (Synthesis~20), 2026. (This archive.)

\bibitem{series:workedexample}
Anonymous.
\newblock Worked Example: Receipt Interlock, Public Packet, and Release-Evolution Stops.
\newblock Series synthesis note (Synthesis~17), 2026. (This archive.)

\bibitem{series:seriesspines}
Anonymous.
\newblock Series Spines from Math to Review: Stable Handoffs for Later Terse Papers.
\newblock Series synthesis note (Synthesis~55), 2026. (This archive.)

\bibitem{series:releaseproperty}
Anonymous.
\newblock Privacy as a Release Property: Proof-Carrying Anonymity Receipts for Anonymous DHTs.
\newblock Release and destination Paper~A, 2026. (This archive.)



% Added in rev0806 citation-closure hygiene pass.
\bibitem{series:eval1}
Anonymous.
\newblock Notarized Anonymity Certificates: Attempt Logs, Spending Discipline, and Verifier Rules.
\newblock Evaluation series Paper~1, 2026. (This archive.)

\end{thebibliography}

\end{document}
